% Options for packages loaded elsewhere
\PassOptionsToPackage{unicode}{hyperref}
\PassOptionsToPackage{hyphens}{url}
\documentclass[
  a4paper,
]{ltjsarticle}
\usepackage{xcolor}
\usepackage[margin=25mm]{geometry}
\usepackage{amsmath,amssymb}
\setcounter{secnumdepth}{-\maxdimen} % remove section numbering
\usepackage{iftex}
\ifPDFTeX
  \usepackage[T1]{fontenc}
  \usepackage[utf8]{inputenc}
  \usepackage{textcomp} % provide euro and other symbols
\else % if luatex or xetex
  \usepackage{unicode-math} % this also loads fontspec
  \defaultfontfeatures{Scale=MatchLowercase}
  \defaultfontfeatures[\rmfamily]{Ligatures=TeX,Scale=1}
\fi
\usepackage{lmodern}
\ifPDFTeX\else
  % xetex/luatex font selection
\fi
% Use upquote if available, for straight quotes in verbatim environments
\IfFileExists{upquote.sty}{\usepackage{upquote}}{}
\IfFileExists{microtype.sty}{% use microtype if available
  \usepackage[]{microtype}
  \UseMicrotypeSet[protrusion]{basicmath} % disable protrusion for tt fonts
}{}
\makeatletter
\@ifundefined{KOMAClassName}{% if non-KOMA class
  \IfFileExists{parskip.sty}{%
    \usepackage{parskip}
  }{% else
    \setlength{\parindent}{0pt}
    \setlength{\parskip}{6pt plus 2pt minus 1pt}}
}{% if KOMA class
  \KOMAoptions{parskip=half}}
\makeatother
\usepackage{longtable,booktabs,array}
\newcounter{none} % for unnumbered tables
\usepackage{calc} % for calculating minipage widths
% Correct order of tables after \paragraph or \subparagraph
\usepackage{etoolbox}
\makeatletter
\patchcmd\longtable{\par}{\if@noskipsec\mbox{}\fi\par}{}{}
\makeatother
% Allow footnotes in longtable head/foot
\IfFileExists{footnotehyper.sty}{\usepackage{footnotehyper}}{\usepackage{footnote}}
\makesavenoteenv{longtable}
\setlength{\emergencystretch}{3em} % prevent overfull lines
\providecommand{\tightlist}{%
  \setlength{\itemsep}{0pt}\setlength{\parskip}{0pt}}
% Glyph map for lualatex (newunicodechar). Maps symbols that may lack font glyphs.
\usepackage{newunicodechar}
\usepackage{amssymb}
\newunicodechar{∎}{\ensuremath{\blacksquare}}
\newunicodechar{✓}{\ensuremath{\checkmark}}
\newunicodechar{✗}{\ensuremath{\times}}
\newunicodechar{⟹}{\ensuremath{\Longrightarrow}}
\newunicodechar{⟺}{\ensuremath{\Longleftrightarrow}}
\newunicodechar{↪}{\ensuremath{\hookrightarrow}}
\newunicodechar{⌊}{\ensuremath{\lfloor}}
\newunicodechar{⌋}{\ensuremath{\rfloor}}
\newunicodechar{≃}{\ensuremath{\simeq}}
\newunicodechar{≈}{\ensuremath{\approx}}
\newunicodechar{≤}{\ensuremath{\leq}}
\newunicodechar{≥}{\ensuremath{\geq}}
\newunicodechar{≠}{\ensuremath{\neq}}
\newunicodechar{→}{\ensuremath{\rightarrow}}
\newunicodechar{←}{\ensuremath{\leftarrow}}
\newunicodechar{↔}{\ensuremath{\leftrightarrow}}
\newunicodechar{⊥}{\ensuremath{\perp}}
\newunicodechar{√}{\ensuremath{\surd}}
\newunicodechar{∑}{\ensuremath{\sum}}
\newunicodechar{∏}{\ensuremath{\prod}}
\newunicodechar{∞}{\ensuremath{\infty}}
\newunicodechar{∈}{\ensuremath{\in}}
\newunicodechar{∀}{\ensuremath{\forall}}
\newunicodechar{∣}{\ensuremath{\mid}}
\newunicodechar{γ}{\ensuremath{\gamma}}
\newunicodechar{λ}{\ensuremath{\lambda}}
\newunicodechar{μ}{\ensuremath{\mu}}
\newunicodechar{ρ}{\ensuremath{\rho}}
\newunicodechar{σ}{\ensuremath{\sigma}}
\newunicodechar{φ}{\ensuremath{\varphi}}
\newunicodechar{ψ}{\ensuremath{\psi}}
\newunicodechar{θ}{\ensuremath{\theta}}
\newunicodechar{Δ}{\ensuremath{\Delta}}
\newunicodechar{Π}{\ensuremath{\Pi}}
\newunicodechar{Λ}{\ensuremath{\Lambda}}
\newunicodechar{τ}{\ensuremath{\tau}}
\newunicodechar{ε}{\ensuremath{\varepsilon}}
\newunicodechar{ω}{\ensuremath{\omega}}
\newunicodechar{π}{\ensuremath{\pi}}
\newunicodechar{ν}{\ensuremath{\nu}}
\newunicodechar{±}{\ensuremath{\pm}}
\newunicodechar{×}{\ensuremath{\times}}
\newunicodechar{·}{\ensuremath{\cdot}}
\usepackage{bookmark}
\IfFileExists{xurl.sty}{\usepackage{xurl}}{} % add URL line breaks if available
\urlstyle{same}
\hypersetup{
  pdftitle={The Additivity Law of the Terminal Equal Amplitude a\_\textbackslash infty=\textbackslash sqrt\{H/M\} ------ Separating Terminal Universality from Ignition Dynamics (N=9 alone vs N=7+N=6{,} 3 initial-value families)},
  pdfauthor={Noriaki Kihara (WF System Co., Ltd.), ORCID 0009-0004-6753-4020},
  hidelinks,
  pdfcreator={LaTeX via pandoc}}

\title{The Additivity Law of the Terminal Equal Amplitude
\(a_\infty=\sqrt{H/M}\) ------ Separating Terminal Universality from
Ignition Dynamics (\(N=9\) alone vs \(N=7+N=6\), 3 initial-value
families)}
\author{Noriaki Kihara (WF System Co., Ltd.), ORCID 0009-0004-6753-4020}
\date{2026-09-12}

\begin{document}
\maketitle

\textbf{Series:} Foundations of Self-Consistent Relational-Wave Closed
Systems (Paper 8 of 10)\\
\textbf{Author:} Noriaki Kihara (WF System Co., Ltd.)　\textbf{Date:}
2026-09-12\\
\textbf{Version DOI:} 10.5281/zenodo.22729116\\
\textbf{Concept DOI:} 10.5281/zenodo.22729115\\
\textbf{ORCID:} 0009-0004-6753-4020\\
\textbf{Zenodo:} https://zenodo.org/records/22729116\\

\begin{center}\rule{0.5\linewidth}{0.5pt}\end{center}

\subsection{Abstract}\label{abstract}

The post-lock terminal state is equal-amplitude (Paper 5). We verify the
\textbf{additivity law} \(a_\infty=\sqrt{H/M}\), in which its radius is
determined solely by the extensive variables \((H,M)\), using the mixed
sum of triangular numbers \(M(9)=M(7)+M(6)\) (\(36=21+15\)). The mixed
system of \(N=7\) and \(N=6\), normalized to the density
\(H_i/M_i=1/36\), agrees in terminal ring radius with the standalone
\(N=9\) system (both giving \(a_\infty=\sqrt{1/36}=1/6=0.1666667\)). On
the other hand, onset, path, and plateau are imprinted by the
composition (\(N=9\) alone onset \(191\) vs mixed union departs at the
earliest stage). Furthermore, changing the initial-value family (staged
\(2\pi/N\) increments / 90-degree analytic floor / 90-degree symmetric
v2) changes \textbf{the presence or absence of ignition itself}
(\(N=9\): staged onset \(191\) / 90-degree analytic floor onset \(140\)
/ 90-degree v2 is the non-igniting exact fixed point \(\sigma=N-1\)),
yet the terminal value \(1/6\) is common across all three families. That
is, \textbf{terminal universality (\(a_\infty=\sqrt{H/M}\)) is
independent of the symmetry family of the initial values}, whereas
\textbf{ignition dynamics depends strongly on it}. The composition
(which combination of \(N\)) is imprinted only during the inflation
period, and at the terminal state it is indistinguishable from a single
system. \((H,M)\) are extensive quantities, \(\rho=H/M\) is an intensive
quantity, and \(a_\infty=\sqrt\rho\) is an intensive terminal state
quantity; the formula itself is an exact consequence of equal amplitude
plus \(H\) conservation, and what is nontrivial is terminal
universality. \textbf{The transient remembers the composition (the
\(J_0,\prod_tJ_t\) of Paper 6), while the terminal state forgets the
composition in its amplitude via \(|z_e|^2\to H/M\), whereas the phase
configuration and the orientation of the face are degenerate and
selected by the seed (Paper 5)} ------ the amplitude becomes a state
quantity, the orientation retains the SSB history. Factorization itself
is dynamically impossible because \(K_N\) is maximally connected
(crossing edges necessarily remain), and the terminal agreement is a
consequence of equipartition unrelated to decomposability.
Classification: numerical verification (additivity law) + mathematical
fact (impossibility of decomposition).

\begin{center}\rule{0.5\linewidth}{0.5pt}\end{center}

\subsection{1. Problem}\label{problem}

If the terminal state is equal-amplitude, then \(a_\infty=\sqrt{H/M}\)
\textbf{follows exactly} from \(H=\sum_{e=1}^M|z_e|^2=Ma_\infty^2\)
------ the formula itself is not nontrivial. The nontrivial part that
this paper confirms by numerical experiment is \textbf{``terminal
universality,'' namely that even different compositions and different
initial-value families reach the same equal-amplitude law at the
terminal state}. Generalization: if the subsystems have identical
density \(\rho_i=H_i/M_i=\rho\), then since
\(H_{\rm tot}=\sum_iH_i=\rho\sum_iM_i=\rho M_{\rm tot}\),

\[\frac{H_{\rm tot}}{M_{\rm tot}}=\rho,\qquad a_{\infty,\rm tot}=\sqrt{H_{\rm tot}/M_{\rm tot}}=\sqrt\rho=a_{\infty,i}.\]

That is, \((H,M)\) are \textbf{extensive quantities}, \(\rho=H/M\) is an
\textbf{intensive quantity}, and \(a_\infty=\sqrt\rho\) is an
\textbf{intensive terminal state quantity}; the additivity of \((H,M)\)
gives the composition-invariance of \(\rho\) (\(N=9=7+6\) is one
instance). We verify this with a heterogeneous mixture and separate
terminal universality from ignition dynamics (the imprint of
composition).

\subsection{2. Dynamics and Design (faithful copy + density
normalization)}\label{dynamics-and-design-faithful-copy-density-normalization}

The map is the canonical one\_step (SHA \texttt{1abf2353…}). The wrapper
is a faithful copy of the canonical implementation with only minimal
changes to the initial-value source and output names (diff audited,
physics formulas unchanged). Using the triangular-number decomposition
\(M(9)=36=21+15=M(7)+M(6)\), the theoretical-floor parents of
\(N=7,N=6\) (each with \(H=1\)) are normalized to the density-matching
condition \(H_i/M_i=1/36\) (\(c_7=\sqrt{21/36}=0.7638\),
\(c_6=\sqrt{15/36}=0.6455\), total \(H=1\)). den=N・500 step. The mixed
system runs each subsystem independently with its own adjacency and its
own den.

\subsection{3. Establishment of the Additivity Law (Experiment
21)}\label{establishment-of-the-additivity-law-experiment-21}

\begin{itemize}
\tightlist
\item
  \textbf{The terminal ring radius agrees}: \(N=9\) alone \(0.1666667\)
  (min=max), the \(N=7\) part of the mixture \(0.1666667\), the \(N=6\)
  part \(0.16649\)--\(0.16677\), all equal to the theoretical value
  \(\sqrt{1/36}=0.1666667\).
\item
  \textbf{The ignition dynamics is entirely different by composition}
  (only the terminal state is universal): onset\((>0.05)\) is \(191\)
  for \(N=9\) alone vs the mixed union departs at the earliest stage
  (the union is meaningless, so it is evaluated per subsystem).
  Per-subsystem onset: \(N=7\) part \(145\), \(N=6\) part \(153\)
  (agreeing with the \(153\) of the \(N=6\times3\) series =
  scale-invariant).
\item
  \(H\) conservation and closure maintenance are sound (\(H\) relative
  difference \(\sim10^{-14}\), maximum closure \(\sim10^{-14}\), finite
  at all steps).
\end{itemize}

Whether by homogeneous \(m\)-fold or heterogeneous mixture, the terminal
state is indistinguishable from a single system via
\(a_\infty=\sqrt{H/M}\). The structure (which combination of \(N\)) is
imprinted only on onset, path, and plateau.

\subsection{4. Change of Ignition Dynamics by Initial-Value Family
(Experiments 23 \& 25) ------ comparison of 3
families}\label{change-of-ignition-dynamics-by-initial-value-family-experiments-23-25-comparison-of-3-families}

The same experiment was re-run changing only the initial-value family:

{\def\LTcaptype{none} % do not increment counter
\begin{longtable}[]{@{}
  >{\raggedright\arraybackslash}p{(\linewidth - 6\tabcolsep) * \real{0.2500}}
  >{\raggedright\arraybackslash}p{(\linewidth - 6\tabcolsep) * \real{0.2500}}
  >{\raggedright\arraybackslash}p{(\linewidth - 6\tabcolsep) * \real{0.2500}}
  >{\raggedright\arraybackslash}p{(\linewidth - 6\tabcolsep) * \real{0.2500}}@{}}
\toprule\noalign{}
\begin{minipage}[b]{\linewidth}\raggedright
\(N=9\)
\end{minipage} & \begin{minipage}[b]{\linewidth}\raggedright
staged (\(2\pi/N\) increments)
\end{minipage} & \begin{minipage}[b]{\linewidth}\raggedright
90-degree analytic floor (integer K)
\end{minipage} & \begin{minipage}[b]{\linewidth}\raggedright
90-degree symmetric v2 (0/90 equal amplitude)
\end{minipage} \\
\midrule\noalign{}
\endhead
\bottomrule\noalign{}
\endlastfoot
ignition & ignition (onset \(191\)) & ignition (onset \(140\)) &
\textbf{non-ignition} (the \(\sigma=N-1\) exact fixed point of
\(N\equiv1\,\mathrm{mod}\,4\)) \\
mixed-part onset (\(N=7,N=6\)) & \(145,\ 153\) & \(113,\ 121\) &
\(5,\ 5\) \\
terminal \(a_\infty\) & \(1/6\) & \(1/6\) & \(1/6\) \\
\end{longtable}
}

\textbf{The terminal value \(1/6\) is common across all three families},
and only the presence/absence of ignition and the onset change by
family. For 90-degree v2, the \(N=9\) case is an exact fixed point where
the \(0^\circ\) family becomes a \((N-1)/2\)-regular graph
(\(\sigma=N-1\)), which sits at rest on the equal-amplitude ring and
does not ignite (consistent with the qualification screening of Paper 2:
a neutral fixed point does not ignite). The 90-degree analytic floor
(Paper 3) and staged ignite on an unstable floor.

\subsection{5. Factorization Is Impossible; the Terminal Agreement Is a
Consequence of Equipartition (Experiment
22)}\label{factorization-is-impossible-the-terminal-agreement-is-a-consequence-of-equipartition-experiment-22}

The complete enumeration of the triangular-number decomposition
\(M(N)=\sum M(k)\) (equal partition, 27 two-part cases, 46 three-part
cases; the ``prime'' \(N\) having no decomposition with parts \(M(4)\)
or larger are \(3,4,5,6,8\)) is the basis for the selection, but
\textbf{dynamical factorization is impossible in principle}: since
\(K_N\) is a complete graph, crossing edges necessarily remain and it is
maximally connected, and \(K_{10}\to3K_6\) violates the degree
condition. The inflation mismatch of Experiment 21 (onset \(191\) vs
\(145/153\)) is the experimental evidence of dynamical impossibility of
decomposition. Therefore, \textbf{the terminal agreement is a
consequence of equipartition (\(a_\infty=\sqrt{H/M}\)) unrelated to
decomposability}. The hierarchization of a condensate = a single vertex
(mass = the \(H\) privately held by the closure one level below) is
presented as a working hypothesis, and its derivation is left as an open
problem for future study.

\subsection{6. Memory and Erasure ------ the Transient Remembers the
Composition and the Terminal Amplitude Forgets It (unification of Papers
5, 6,
8)}\label{memory-and-erasure-the-transient-remembers-the-composition-and-the-terminal-amplitude-forgets-it-unification-of-papers-5-6-8}

The carrier of information separates between the transient and the
terminal state. The \textbf{transient} (inflation period) remembers the
composition and initial-value family: the initial structure is imprinted
in the floor Jacobian \(J_0\) and the tangent cocycle \(\prod_t J(z_t)\)
(Paper 6), determining onset, path, and plateau (§3--4: \(N=9\) alone
onset \(191\) vs mixed \(145/153\), and even the presence/absence of
ignition changes with staged/analytic floor/v2). The \textbf{terminal
state} has \(|z_e|^2\to H/M\) after nonlinear saturation, and the
composition information disappears from this amplitude
(\(a_\infty=\sqrt{H/M}\) is common across 3 families × mixture). That
is,

\[\text{initial structure}\ \to\ \text{memorized in}\ J_0,\ \textstyle\prod_t J_t\ \to\ \text{inflation}\ \to\ \text{saturation}\ \to\ \text{universalization of amplitude information}.\]

Connection with Paper 5 (amplitude forgets, orientation remains): the
degenerate vacuum family of Paper 5 possesses both a \textbf{common
amplitude scale} \(|z_e|=\sqrt{H/M}\) and a \textbf{seed-dependent phase
configuration and orientation of the face (geometric moduli)}. The
radial coordinate \(r=\sqrt{H_\perp/H}\) of Paper 5 (\(0.19\)--\(0.38\)
per vacuum) is the order parameter of the tilt of the face, and it is a
\textbf{distinct quantity} from the per-edge amplitude
\(a_\infty=\sqrt{H/M}\) of this paper (they are not identified).
Combining the two, the terminal appearance is

\[\boxed{\text{the amplitude scale }|z_e|=\sqrt{H/M}\text{ is universal (becomes a state quantity) / the phase configuration and the orientation of the face are degenerate and selected by the seed (retain the SSB history)}}\]

------ \textbf{at the terminal state the amplitude forgets the
composition, and the orientation retains the history of SSB}. This is
the unification of Paper 5 (degenerate vacuum), Paper 6 (transient
inflation), and Paper 8 (terminal state quantity).

\subsection{7. Claim Classification \& Open
Items}\label{claim-classification-open-items}

\begin{itemize}
\tightlist
\item
  Classification: numerical verification (additivity law
  \(a_\infty=\sqrt{H/M}\), 3 families × mixture) + mathematical fact
  (impossibility of factorization). Verdict: retained.
\item
  Note: the terminal equal amplitude has the phase normalization of this
  engine (stage 2, which kills the amplitude) built in, so it is not
  presented as an independent corroboration of the norm form (suspicion
  of circular derivation; independent verification with an
  amplitude-preserving control dynamics is open).
\item
  Verified (re-run C): the 3-way decomposition of the edge space
  \(1\oplus(N-1)\oplus N(N-3)/2\) (line graph / Johnson scheme) exists
  in reality and the terminal \(\sigma_{\max}\to N-1\) is confirmed for
  all \(N\). However, the terminal state does not ride on the
  \((N-1)\)-dimensional vertex sector but spreads into the residual
  sector (frac\_rest \(\to0.93\)) ------ \textbf{``\(\sigma=N-1\) (the
  eigenvalue) ↔ vertex sector (the dimension)'' is a coincidence of
  dimension labels and does not correspond} (rejected).
\item
  Open: quantification of the immediate departure of the mixed system
  (union onset).
\end{itemize}

\subsection{Related work (structural agreement, not the source of
derivation)}\label{related-work-structural-agreement-not-the-source-of-derivation}

\begin{itemize}
\tightlist
\item
  Barycentric decomposition of mean-field / all-pairs-coupled systems
  (\(K_k\) spectrum \(\{0,k,\ldots,k\}\)): the reduction N-body =
  conserved barycenter + independent two-body problem.
\item
  Extensive / intensive variables (thermodynamics):
  \(a_\infty=\sqrt{H/M}\) is a state quantity of the extensive variables
  \((H,M)\).
\end{itemize}

\subsection{References}\label{references}

\begin{enumerate}
\def\labelenumi{\arabic{enumi}.}
\tightlist
\item
  Noriaki Kihara, ``The Mechanism of Inflationary Rapid Expansion in
  Self-Consistent Relational-Wave Closed Systems,'' Concept DOI
  10.5281/zenodo.22112008.
\item
  Noriaki Kihara, ``Onset-Mode Discrimination (Amplification and
  Unstable Relative Equilibria),'' Concept DOI 10.5281/zenodo.21798854.
\end{enumerate}

\subsection{Reproduction}\label{reproduction}

Wrappers, parents, and figures are packaged together in each set:
\texttt{ロールバック対照テスト\_正本実装N3\_N6\_20260908/} (Experiments
20 \& 21: \texttt{run\_single\_simplex\_N9\_v1.py},
\texttt{make\_normalized\_parents\_N7N6\_v1.py},
\texttt{run\_mixed\_simplex\_N7N6\_v1.py},
\texttt{plot\_Hperp\_mixed\_parts\_N7N6\_v1.py}, 8 figures),
\texttt{加法法則\_90度系\_N9\_vs\_N7N6\_20260909/} (Experiment 23),
\texttt{加法法則\_90度解析床\_N9\_vs\_N7N6\_20260909/} (Experiment 25),
\texttt{二体三体問題\_関係波の検討\_20260908/} (Experiment 22:
\texttt{enumerate\_triangular\_decompositions\_v1.py},
\texttt{三角数分解一覧\_N3\_N40.txt},
\texttt{分析\_二体三体問題\_関係波と凝縮体頂点階層\_20260908.md}). The
test of the correspondence between the vertex sector and \(\sigma=N-1\)
(re-run C) is in
\texttt{位相ロック全N確認\_相対平衡\_20260912/頂点セクター\_sigmaN1\_20260912/}
(\texttt{analyze\_vertex\_sector\_sigmaN1\_20260912.py},
\texttt{results/vertex\_sector\_summary.csv}, README / REPORT / SHA /
run\_all; read-only). Each includes README, SHA256SUMS, run\_all.sh, and
execution logs. The wrapper is a faithful copy of the canonical
\texttt{run\_N3\_N40\_stage123\_v1.py} (SHA256
\texttt{1abf2353fee2e4f56f05e7a6f149fd086885136beb61ab571b48a56b09691567})
with only minimal changes to the initial values and output names.

\end{document}
